\section{Conclusion}

We presented a deep-learning pipeline for variable-length N-CMAPSS turbofan RUL prediction and used it to compare recurrent, convolutional, hybrid, transformer, and multi-scale attention models under a common preprocessing, training, and evaluation harness. The final review-response suite is a simulation study on FD1 with three seeds, reader max RUL 65, and explicit documentation of virtual-sensor use. WaveNet is the strongest controlled FD1 model tested here (RMSE $6.868\pm0.299$ at $T{=}1000$), MSTCN is stable and competitive (RMSE $7.890\pm0.526$), and CNN-GRU is not reliable because two seeds collapse. A matched WaveNet length sweep shows that short windows are effective across 100--1000 timesteps, with $T{=}500$ best by mean RMSE ($6.690\pm0.173$). Asymmetric MSE is modestly better than symmetric MSE in this limited suite, but broader loss and architecture ablations remain open. The code, configs, benchmark outputs, and W\&B traces provide a reproducible basis for extending this comparison; the results should not be read as direct evidence for real-world commercial-engine deployment or as a state-of-the-art claim against studies using different protocols.

\subsection*{Limitations}

\begin{itemize}
  \item The sequence-length result is now matched for 100, 250, 500, and 1000 timesteps, but still needs 2000, 5000, and full-flight inputs under the same architecture.
  \item MSTCN component attribution is incomplete; the multi-scale branches and global fusion attention should be removed independently.
  \item The asymmetric-loss result is limited to WaveNet on FD1; it needs an $\alpha$ sensitivity sweep and replication across architectures and subsets.
  \item Evaluation is currently concentrated on FD1 and FD2; FD3--FD7 remain necessary for broad N-CMAPSS generalization claims.
  \item Comparisons to STAR, TTSNet, and Xu et al.'s MSTCN are architectural rather than apples-to-apples because those papers report primarily on C-MAPSS, not N-CMAPSS.
\end{itemize}
